\setcounter{equation}{0}
\setcounter{figure}{0}
\setcounter{table}{0}
\setcounter{page}{1}
\makeatletter
\renewcommand{\thesection}{S\arabic{section}}
\renewcommand{\theequation}{S\arabic{equation}}
\renewcommand{\thefigure}{S\arabic{figure}}

\begin{widetext}
\begin{center}
\textbf{\large Supplemental Material: Mapping out phase diagrams with generative classifiers}
% \\~\\
% Julian Arnold$^{1}$, Frank Sch\"afer$^{2}$, Alan Edelman$^{2,3}$, and Christoph Bruder$^{3}$ \\
% \vspace{.05in}
% \small{
% $^{1}$\textit{Department of Physics, University of Basel, Klingelbergstrasse 82, 4056 Basel, Switzerland}\\
% $^{2}$\textit{CSAIL, Massachusetts Institute of Technology, Cambridge, MA
% 	02139, USA}\\
% $^{3}$\textit{Department of Mathematics, Massachusetts Institute of Technology, Cambridge, MA
% 	02139, USA}
% }
\end{center}
\end{widetext}

\section{Additional details on phase-classification methods}
\subsection{Class imbalance}
In traditional discriminative machine learning (ML) tasks, class imbalance occurs when the number of samples representing each class within the data set differs. In phase-classification tasks, this can, for example, occur when one does not have access to an equal amount of data at each point in parameter space. Moreover, in scheme 1 (with indicator $I_{1}$) and 2 (with indicator $I_{2}$) discussed in the main text, an imbalance can also occur when sampling an \emph{equal} amount of data at each point in parameter space. In both cases, the imbalance arises because the sets $\{ \Gamma_{y}\}_{y \in \mathcal{Y}}$ are not of equal size. In scheme 1, this can be circumvented by choosing an equal number of points in parameter space to represent each phase. In scheme 2, data sets of unequal size occur as the bipartitions artificially divide the uniformly sampled one-dimensional parameter space. In previous works~\cite{van:2017,liu:2018,beach:2018,suchsland:2018,lee:2019,greplova:2020,bohrdt:2021,richter:2022,gavreev:2022,zvyagintseva:2022,schlomer:2023}, this class imbalance was not accounted for. As a result, the corresponding indicator $I_{2}$ will exhibit trivial local maxima ($I_{2}=1$) at the edges of the sampled region of the parameter space where the entire data is labeled 1 or 2, respectively. If no phase transition is present, the indicator will thus exhibit a characteristic V-shape. In the presence of a phase transition, a W-shape is expected, where the location of the intermediate maximum corresponds to the predicted critical point. Such behavior is undesirable for an indicator of phase transitions, and can even lead to a failure to detect certain transitions (or misdetection), see Fig.~\ref{fig:LBC_imbalance} for an example. In Ref.~\cite{bohrdt:2021}, an attempt has been made to correct this pathological behavior. However, \citet{arnold:2022} have found that this procedure biases the transition point towards the center of the parameter range under consideration.

 \begin{figure}[bth!]
	\centering
		\includegraphics[width=0.9\linewidth]{fig_LBC_bias.pdf}
		\caption{Results for the cluster-Ising model [Eq.~(9) in the main text, $L=7$] at $h_{1}/J = 0.2$ measured using a Pauli-6 POVM. The set $\Gamma$ is composed of a uniform grid with 101 points. Indicator $I_{2}(\gamma)$ with $l = 101$ (i.e., considering global bipartitions of the one-dimensional parameter space) and in the presence (blue) or absence (black) of class imbalance. Estimated critical points (red dashed lines) are determined from maxima in $|\partial^2 \langle H \rangle_{\bm{\gamma}} /\partial \gamma_{2}^2|$. Here, the expected value involved in $I_{2}(\gamma)$ is computed exactly and the ground state is obtained via exact diagonalization.}
	\label{fig:LBC_imbalance}
\end{figure}

In this work, we addressed class imbalances arising from sets $\{ \Gamma_{y}\}_{y \in \mathcal{Y}}$ of unequal size by choosing a uniform prior distribution $P(y)$. In particular, this removes the trivial local maxima at the edges of the sampled region of the parameter space previously encountered for the indicator $I_{2}$ and we instead expect a single peak in the presence of a phase transition. Figure~\ref{fig:LBC_imbalance} shows an example of transitions that are not detected using scheme 2 in the presence of class imbalance highlighting the importance of our proposed modification.

\subsection{Scheme 3: Dividing by the standard deviation}
For the third phase-classification scheme discussed in the main text (with indicator $I_{3}$), we have proposed to divide the corresponding signal by its standard deviation. We find this modified indicator to be superior compared to its unmodified version
\begin{align*}\label{eq:IPBM_SI}
    I_{3}'(\bm{\gamma}) &= \sqrt{\sum_{i=1}^{d} \left(\frac{\partial \hat{\gamma}_{i}(\bm{\gamma})}{\partial \gamma_{i}}\right)^2},\\
    &=  \norm{ \mathbb{E}_{\bm{x} \sim P(\bm{x}|\bm{\gamma})}\left[\hat{\bm{\gamma}}(\bm{x}) \nabla_{\bm{\gamma}} \ln P(\bm{x}|\bm{\gamma}) \right]}_{2} \numberthis,
\end{align*}
which (up to a constant offset) has been considered in previous studies~\cite{schaefer:2019,greplova:2020,arnold:2021,arnold:2022}. As an example, Fig.~\ref{fig:fig_PBM_mod} compares the two indicators for the anisotropic Ising model. The indicator $I_{3}'$ shows four distinct peaks with the dominant two occurring within the two ordered phases (i.e., at larger $|J_{y}/k_{\rm B}T|$). Such spurious signals have also been observed in previous studies using indicators where the standard deviation is not taken into account~\cite{schaefer:2019,arnold:2022}. In contrast, the modified indicator $I_{3}$ we propose in this work shows two distinct peaks which qualitatively agree with the two critical points.

 \begin{figure}[tbh!]
	\centering
		\includegraphics[width=0.99\linewidth]{fig_PBM_mod.pdf}
		\caption{Results for the anisotropic Ising model on a square lattice ($L=20$) at $J_{x}/k_{\rm B}T = -0.325$. The indicator $I_{3}(\gamma)$ (black, left $y$-axis) and $I_{3}'(\gamma)$ (blue, right $y$-axis) are computed based on predictions $\hat{\bm{\gamma}}$ that are estimated element-wise from line scans across the entire two-dimensional phase diagram. The set $\Gamma$ is composed of a uniform grid with 60 points for each axis. Each data set $\mathcal{D}_{\bm{\gamma}}$ consists of $10^5$ spin configurations. Onsager's analytical solution for the phase boundaries is shown as a red dashed line.}
	\label{fig:fig_PBM_mod}
\end{figure}

\subsection{Discriminative approach to compute indicators}\label{sec:discr_indicators}
The computation of all three indicators can be approached by solving classification tasks. These tasks can be tackled in a discriminative manner. In this case, we look for the parameters $\bm{\theta}$ of a parametric model $\tilde{P}_{\bm{\theta}}(y|\bm{x})$ that minimize the following cross-entropy loss function
\begin{equation}\label{eq:SI_x1}
    \mathcal{L}(\bm{\theta}) = - \frac{1}{|\mathcal{Y}|} \sum_{y \in \mathcal{Y}} \frac{1}{|\mathcal{D}_{y}|} \sum_{\bm{x} \in \mathcal{D}_{y}}  {\rm ln}\left(\tilde{P}_{\bm{\theta}}(y|\bm{x})\right),
\end{equation}
where $\mathcal{D}_{y} = \{ \bm{x} \in \mathcal{D}_{\bm{\gamma}} |\bm{\gamma} \in \Gamma_{y}\}$ is the relevant data set drawn from $P(\bm{x}|y)=1/|\Gamma_{y}| \sum_{\bm{\gamma} \in \Gamma_{y}} P(\bm{x}|\bm{\gamma})$ and $\mathcal{D}_{\bm{\gamma}}$ corresponds to a set of samples drawn from $P(\bm{x}|\bm{\gamma})$. Here, class imbalance is addressed by the rescaling factors $1/|\mathcal{D}_{y}|$ as is common in neural-network training~\cite{paszke:2019}. 

Based on the trained parametric model, in scheme 1 we can estimate the central quantity of interest $P(y|\bm{\gamma})$ as 
\begin{equation}\label{eq:SL_1}
    P(y|\bm{\gamma}) \approx \frac{1}{|\mathcal{D}_{\bm{\gamma}}|}\sum_{\bm{x} \in \mathcal{D}_{\bm{\gamma}}} \tilde{P}_{\bm{\theta}}(y|\bm{x})
\end{equation}
In scheme 2, the relevant quantity $p_{\rm err}$ can, for example, be estimated as
\begin{equation}\label{eq:LBC_err_SI}
    p_{\rm err} \approx \frac{1}{2}\sum_{y \in \{1,2 \}} \frac{1}{|\mathcal{D}_{y}|}\sum_{\bm{x} \in \mathcal{D}_{y}} \left|y- \argmax_{y'} \tilde{P}_{\bm{\theta}}(y'|\bm{x})\right|.
\end{equation}
In scheme 3, the key quantity is $\hat{\bm{\gamma}}(\bm{x})$ which can be estimated as
\begin{equation}\label{eq:PBM_pred_SI}
    \hat{\bm{\gamma}}(\bm{x}) \approx \sum_{y \in \mathcal{Y}} \tilde{P}_{\bm{\theta}}(y|\bm{x})\bm{\gamma}_{y}
\end{equation}
Alternatively, it may also be approximated directly using a parametric predictive model $\hat{\bm{\gamma}}_{\bm{\theta}}(\bm{x})$ that is trained to solve a regression task instead of a classification task (as has been done in previous works~\cite{schaefer:2019,greplova:2020,arnold:2021,arnold:2022}). In this case, the relevant loss function is
\begin{equation}\label{eq:SI_x2}
    \mathcal{L}(\bm{\theta}) = \frac{1}{|\mathcal{Y}|} \sum_{y \in \mathcal{Y}} \frac{1}{|\mathcal{D}_{y}|} \sum_{\bm{x} \in \mathcal{D}_{y}}  \norm{\hat{\bm{\gamma}}_{\bm{\theta}}(\bm{x}) - \bm{\gamma}_{y}}_{2}^2.
\end{equation}

Having introduced loss functions that can be used to optimize parametric predictive models, in the following we prove that in the infinite-data limit, an optimal predictive model that minimizes the corresponding loss functions results in the indicator signals discussed in the main text.

\indent For a given sample $\bm{x} \in \{x \in \mathcal{D}_{y}|y \in \mathcal{Y} \}$, we can determine the corresponding empirically optimal model prediction $P_{\rm emp}(y|\bm{x})$ by minimizing the loss function in Eq.~\eqref{eq:SI_x1} with respect to $P_{\rm emp}(y|\bm{x})$ subjected to the equality constraint $\sum_{y\in\mathcal{Y}} P_{\rm emp}(y|\bm{x}) - 1 = 0$. Using the method of Lagrange multipliers, the stationary points satisfy the following conditions
\begin{equation}
    \lambda = \frac{1}{|\mathcal{Y}|} \frac{M_{y}(\bm{x})}{|\mathcal{D}_{y}|} \frac{1}{P_{\rm emp}(y|\bm{x})}\; \forall y \in \mathcal{Y},
\end{equation}
where $M_{y}(\bm{x})$ denotes the number of times $\bm{x}$ appears in $\mathcal{D}_{y}$ and $\lambda$ is the corresponding Lagrange multiplier. Together with the equality constraint, this yields
\begin{equation}
    P_{\rm emp}(y|\bm{x}) = \frac{\frac{M_{y}(\bm{x})}{|\mathcal{D}_{y}|} \frac{1}{|\mathcal{Y}|}}{\sum_{y' \in \mathcal{Y}} \frac{M_{y'}(\bm{x})}{|\mathcal{D}_{y'}|} \frac{1}{|\mathcal{Y}|}}.
\end{equation}
In the infinite-data limit $|\mathcal{D}_{y}| \rightarrow \infty$, 
 we have $M_{y}(\bm{x})/|\mathcal{D}_{y}| \rightarrow P(\bm{x}|y)$. Thus, the predictions of the empirically optimal model converge to the predictions of a Bayes optimal model
\begin{equation}\label{eq:SI_optt}
    P_{\rm emp}(y|\bm{x}) \rightarrow P(y|\bm{x}) = \frac{P(\bm{x}|y)P(y)}{\sum_{y' \in \mathcal{Y}} P(\bm{x}|y')P(y')},
\end{equation}
with a uniform prior over the classes $P(y) = 1/|\mathcal{Y}|\; \forall y \in \mathcal{Y}$. Thus, we have recovered Eq.~(1) of the main text.

\indent Following a similar procedure for the loss function in Eq.~\eqref{eq:SI_x2}, we obtain
\begin{equation}
    \hat{\bm{\gamma}}_{\rm emp}(\bm{x}) = \sum_{y \in \mathcal{Y}} \frac{ \frac{M_{y}(\bm{x})}{|\mathcal{D}_{y}|} \frac{1}{|\mathcal{Y}|}}{\sum_{y' \in \mathcal{Y}} \frac{M_{y'}(\bm{x})}{|\mathcal{D}_{y'}|} \frac{1}{|\mathcal{Y}|}} \bm{\gamma}_{y},
\end{equation}
which in the infinite-data limit converges to 
\begin{equation}
    \hat{\bm{\gamma}}(\bm{x}) = \sum_{y\in \mathcal{Y}} \frac{P(\bm{x}'|y)P(y)}{\sum_{y'\in \mathcal{Y}} P(\bm{x}'|y')P(y')} \bm{\gamma}_{y} = \sum_{y\in \mathcal{Y}} P(y|\bm{x}) \bm{\gamma}_{y},
\end{equation}
corresponding to Eq.~(6) in the main text.

\subsection{Discriminative cooperative networks}
\indent In Ref.~\cite{liu:2018}, a scheme for mapping out phase diagrams in two-dimensional parameter spaces was proposed. Instead of a brute-force search of the entire parameter space for all possible phase boundaries, the predicted phase boundary is modeled as a parametrized curve using an active contour model called \emph{snake}~\cite{kass:1988}. This snake partitions the parameter space locally and is driven via internal forces that, e.g., prevent bending and stretching, as well as external forces aiming to minimize the overall classification error. In Ref.~\cite{liu:2018}, the external force is generated in an interplay between a guesser network and learner network, together called \emph{discriminative cooperative networks}, that are optimized via a joint cost function. The guesser provides labels for the data which the learner should reproduce. In each step, the guesser tries to provide a better set of labels based on the predictions of the learner to cooperatively minimize the classification error (i.e., the corresponding loss).

\indent Let us first consider a one-dimensional parameter space. We use the following sigmoidal guesser
\begin{equation}
    \tilde{P}_{\bm{\theta}_{G}}(1|\gamma) = \frac{1}{1+e^{(\gamma_{G} - \gamma)/\sigma_{G}}},
\end{equation}
where $\tilde{P}_{\bm{\theta}_{G}}(2|\gamma) = 1 - \tilde{P}_{\bm{\theta}_{G}}(1|\gamma)$. The guesser is characterized by two parameters $\bm{\theta}_{G}=(\gamma_{G},\sigma_{G})$, where $\gamma_{G}$ corresponds to the guessed transition point and $\sigma_{G}$ determines the sharpness of the transition. The learner is given by $\tilde{P}_{\bm{\theta}_{L}}(y|\bm{x})$. The guesser and learner are optimized jointly using a cross-entropy loss function
\begin{widetext}
\begin{equation}\label{eq:DCN}
    \mathcal{L}(\bm{\theta}) = -  \frac{1}{|\mathcal{Y}|}\sum_{y \in \mathcal{Y}} \frac{1}{|\Gamma_{y}|}\sum_{\gamma \in \Gamma_{y}}\frac{1}{|\mathcal{D}_{\gamma}|}\sum_{\bm{x} \in \mathcal{D}_{\gamma}} \tilde{P}_{\bm{\theta}_{G}}(y|\gamma) \log\left(\tilde{P}_{\bm{\theta}_{L}}(y|\bm{x})\right),
\end{equation}
\end{widetext}
where $\bm{\theta} = \left( {\bm{\theta}_{G}}, {\bm{\theta}_{L}}\right)$. The parameters of the two networks can now be optimized using gradient descent on the loss function in Eq.~\eqref{eq:DCN}. This corresponds to a discriminative approach and has been used in Ref.~\cite{liu:2018}.

\indent Moving to two-dimensional parameter spaces, one can utilize the snake as a parametrization for the guesser which is a discretized curve of linked nodes, $\bm{r}(s)=\left(\gamma_{1}(s),\gamma_{2}(s)\right)$, parametrized by $s \in [0,1]$ (assuming an open snake), see Fig.~\ref{fig:DCNs_LBC}(a). The snake moves in order to minimize its total energy $E_{\rm tot}=E_{\rm int} + E_{\rm ext}$. The internal energy
\begin{equation}
    E_{\rm int} = \int_0^1\left(\alpha \norm{\frac{\partial \bm{r}}{\partial s}}_{2}^2+\beta\norm{\frac{\partial^2 \bm{r}}{\partial s^2}}_{2}^2\right) ds
\end{equation}
is introduced to make the snake smoother, where the hyperparameter $\alpha$ penalizes stretching of the snake and $\beta$ penalizes bending. The snake can sense its surrounding at each node through $2l$ sampled points perpendicular to the snake within a distance $2l \sigma$ [see Fig.~\ref{fig:DCNs_LBC}(b)]. The overall guesser function is then comprised by local guesser functions evaluated at each node that senses the direction perpendicular to the snake. The external energy $E_{\rm ext}$ of the snake corresponds to the overall loss obtained by summing the cross-entropy losses of all individual one-dimensional guessers [cf. Eq.~\eqref{eq:DCN}] and gives rise to an external force $-\delta E_{\rm ext}/\delta \mathbf{r}$ pointing perpendicular to the snake at each node.

\indent In this work, we replace the learner network at each node by a corresponding generative classifier which makes updating the learner in each step obsolete. To construct the generative classifier, we derive the (empirically) optimal predictor based on the loss in Eq.~\eqref{eq:DCN}. Following the same procedure as in the previous section, we obtain 
\begin{equation}
    P_{\rm emp}(y|\bm{x}) = \frac{\sum_{\gamma \in \Gamma_{y}} \tilde{P}_{\bm{\theta}_{G}}(y|\gamma) \frac{M_{\gamma}(\bm{x})}{|\mathcal{D}_{\gamma}|} \frac{1}{|\mathcal{Y}|}\frac{1}{|\Gamma_{y}|}}{\sum_{y' \in \mathcal{Y}} \sum_{\gamma' \in \Gamma_{y'}} \tilde{P}_{\bm{\theta}_{G}}(y'|\gamma') \frac{M_{\gamma'}(\bm{x})}{|\mathcal{D}_{\bm{\gamma}'}|} \frac{1}{|\mathcal{Y}|}\frac{1}{|\Gamma_{y'}|}}.
\end{equation}
In the infinite-data limit, this converges to 
\begin{equation}\label{eq:DCN_gen}
    P(y|\bm{x}) = \frac{\sum_{\gamma \in \Gamma_{y}} \tilde{P}_{\bm{\theta}_{G}}(y|\gamma)P(\bm{x}|\gamma)P(y)P(\gamma|y)}{\sum_{y'\in \mathcal{Y}} \sum_{\gamma' \in \Gamma_{y'}}\tilde{P}_{\bm{\theta}_{G}}(y'|\gamma')P(\bm{x}|\gamma')P(y)P(\gamma'|y')},
\end{equation}
with a uniform prior over the classes $P(y) = 1/|\mathcal{Y}|\; \forall y \in \mathcal{Y}$ and $P(\gamma|y) = 1/|\Gamma_{y}|$ if $\gamma \in \Gamma_{y}$ and zero otherwise. Thus, we can obtain a generative classifier by modeling $P(\bm{x}|\gamma)$ in Eq.~\eqref{eq:DCN_gen}. Note that the expression in Eq.~\eqref{eq:SI_optt} for binary labels is recovered as a special case.\\

\begin{figure*}[htb!]
	\begin{center}
		\includegraphics[width=0.9\linewidth]{fig_SI_DCN.png}
		\caption{(a) Schematic illustration of the snake model. The blue circles represent the snake nodes (here 7). The normal (sensing) direction at each node is shown as a black-dashed line and the relevant $2l$ sampled points within a distance $2l \sigma$ are shown as black circles (here $l=2$). (b) Snake boundary for the anisotropic Ising model on a square lattice ($L=20$; $J_{x}/k_{\rm B}T,J_{y}/k_{\rm B}T\geq 0$) obtained using a generative classifier where the color bar denotes the training epoch. The snake consists of 20 nodes with $l=4$ and hyperparameters $\alpha = 0.002$, $\beta = 0.4$, $\kappa = 0.9$, and learning rates of $10^{-4}$ and $5 \times 10^{-4}$ associated with $E_{\rm int}$ and $E_{\rm ext}$, respectively. The set $\Gamma$ is composed of a uniform grid with 30 points for each axis. The data sets $\mathcal{D}_{\bm{\gamma}}$ consist of $10^5$ spin configurations each. Onsager's analytical solution for the phase boundary is shown as a black-dashed line.
  }
		\label{fig:DCNs_LBC}
	\end{center}
\end{figure*}

In our implementation, we fix the width parameter $\sigma_{G}$ of the sigmoid guesser to be $\sigma_{G} = \sigma/10$, where $\sigma$ is the width parameter of the snake. We reduce $\sigma$ exponentially during the optimization according to 
\begin{equation}
    \sigma_{k} = \sigma_{\rm end}+(\sigma_{\rm start}-\sigma_{\rm end}) \kappa^{k},
\end{equation}
where $k$ denotes the current epoch, $\sigma_{\rm start/end}$ are the start and end values, and $\kappa$ is the decay rate. We set $\sigma_{\rm end}=\Delta \gamma$, where $\Delta \gamma$ is the spacing between neighboring sampled points in parameter space and $\sigma_{\rm start} = 5 \sigma_{\rm end}$. This choice yields a strong gradient signal early on that becomes weaker but more accurate at later stages, which facilitates the convergence of the snake to the true underlying phase boundary. We construct models for $P(\bm{x}|\bm{\gamma})$ at points $\bm{\gamma}$ lying off the initial sampled grid using bilinear interpolation. Gradients of the energy with respect to the node positions are calculated using finite differencing. We minimize the snake's total energy using gradient-based optimization with Adam~\cite{kingma:2014}.\\ 

The results obtained using this scheme with a generative learner for the anisotropic Ising model are shown in Fig.~\ref{fig:DCNs_LBC}(b). While the snake eventually finds the underlying phase boundary, there are several downsides to this scheme. First, the number of unique classification tasks that need to be solved over the course of the training is $N_{\rm epochs} \times N_{\rm nodes} = 400 \times 20 = 8000$. This is larger than the number of unique classification tasks that need to be solved within scheme 3 proposed in the main text which is given by $|\Gamma| \times 2 = 900 \times 2 = 1800$. Therefore, there is no gain in computation time compared to a brute-force search of the parameter space (as performed within scheme 3). The snake scheme also has a multitude of hyperparameters: $\alpha$, $\beta$, $\sigma_{\rm start}$, $\kappa$, $l$, learning rates, as well as the initial snake positioning and snake topology (e.g., whether the snake is open or closed and whether certain nodes are being fixed). Thus, to run the snake scheme one has to perform hyperparameter tuning. This involves additional computational effort and/or prior knowledge of the underlying phase diagram. In particular, we find the results to depend heavily on $\alpha$, $\beta$, as well as the learning rates and initial position of the snake. Moreover, the scheme is not reliable in the presence of more than two phases as the snake will typically converge to one of the phase boundaries [see Fig.~\ref{fig:DCNs_LBC}(b)]. To get around this, the scheme must be run multiple times with different snake initialization (possibly guided by prior knowledge of the phase diagram).\\  

\section{Mapping out phase diagrams of classical equilibrium systems}
In this section, we prove that to map out phase diagrams of a large class of classical equilibrium systems by computing one of the three indicators discussed in the main text, it suffices to model $\{ P(\bm{X}|\bm{\gamma})\}_{\bm{\gamma} \in \Gamma}$, where $\bm{X}$ is the corresponding sufficient statistic.

Consider probability distributions of the exponential family
\begin{equation}\label{eq:SM_exp_fam}
    P_{\rm exp}(\bm{x}|\bm{\gamma}) = h(\bm{x})\exp(\sum_{i=1}^{d} \eta_{i}(\bm{\gamma}) X_{i}(\bm{x})-A(\bm{\gamma})),
\end{equation}
where $h(\bm{x})\geq 0$ is the carrier measure, $\bm{X}(\bm{x})$ is a sufficient statistic, $\bm{\eta}(\bm{\gamma})$ are the natural parameters, and $A(\bm{\gamma})$ is the log-partition function. The statistic $\bm{X}(\bm{x})$ is a \emph{minimal} sufficient statistic for $\bm{\gamma}$ if the set of allowed natural parameters $\bm{\eta}(\bm{\gamma})$ spans a $d$-dimensional space~\cite{casella:2002}. The distribution over the sufficient statistic is given by
\begin{align*}
    P_{\rm exp}(\bm{X}'|\bm{\gamma}) &= \sum_{\bm{x} \in \mathcal{X}} P_{\rm exp}(\bm{x}|\bm{\gamma}) \delta(\bm{X}(\bm{x}) - \bm{X}')\\
    &= g(\bm{X}')\exp(\bm{\eta}(\bm{\gamma}) \cdot \bm{X}' - A(\bm{\gamma}))\numberthis,
\end{align*}
where $g(\bm{X}') = \sum_{\bm{x} \in \mathcal{X}} h(\bm{x})\delta(\bm{X}(\bm{x}) - \bm{X}')$. We have 
\begin{equation}\label{eq:SM_der}
    P_{\rm exp}(\bm{x}|\bm{\gamma})/P(\bm{X}(\bm{x})|\bm{\gamma})=h(\bm{x})/g(\bm{X}(\bm{x})).
\end{equation}
\begin{widetext}
Crucially, this ratio is independent of $\bm{\gamma}$. Using Eq.~\eqref{eq:SM_der}, the conditional probability of a label in phase-classification tasks [Eq.~(2) in the main text] can be written as
\begin{equation}\label{eq:Bayes_2_SM}
    P(y|\bm{x}) =  \frac{\frac{1}{|\Gamma_{y}|}\sum_{\bm{\gamma} \in \Gamma_{y}} P_{\rm exp}(\bm{x}|\bm{\gamma})}{\sum_{y' \in \mathcal{Y}}\frac{1}{|\Gamma_{y'}|}\sum_{\bm{\gamma}' \in \Gamma_{y'}} P_{\rm exp}(\bm{x}|\bm{\gamma}')}
     =  \frac{\frac{1}{|\Gamma_{y}|}\sum_{\bm{\gamma} \in \Gamma_{y}} P(\bm{X}(\bm{x})|\bm{\gamma})}{\sum_{y' \in \mathcal{Y}}\frac{1}{|\Gamma_{y'}|}\sum_{\bm{\gamma}' \in \Gamma_{y'}} P(\bm{X}(\bm{x})|\bm{\gamma}')}
    = P(y|\bm{X}(\bm{x})).
\end{equation}
Thus,
\begin{align*}\label{eq:SM_der_1}
    P(y|\bm{\gamma}) &= \mathbb{E}_{\bm{x} \sim P_{\rm exp}(\bm{x}|\bm{\gamma})}\left[ P(y|\bm{x}) \right] =  \sum_{\bm{x} \in \mathcal{X}} P_{\rm exp}(\bm{x}|\bm{\gamma})  P(y|\bm{x}) = \sum_{\bm{x} \in \mathcal{X}}  P(\bm{X}(\bm{x})|\bm{\gamma})  P(y|\bm{X}(\bm{x})) h(\bm{x})/g(\bm{X}(\bm{x}))\\
    &= \sum_{\bm{X}' \in \mathcal{X}_{\rm suff}}\sum_{\bm{x} \in \mathcal{X}} P(\bm{X}'|\bm{\gamma})  P(y|\bm{X}') \delta(\bm{X}(\bm{x})-\bm{X}') h(\bm{x})/g(\bm{X}')= \sum_{\bm{X} \in \mathcal{X}_{\rm suff}} P(\bm{X}|\bm{\gamma})  P(y|\bm{X}) = \mathbb{E}_{\bm{X} \sim P(\bm{X}|\bm{\gamma})}\left[  P(y|\bm{X})\right] \numberthis, \\
\end{align*}
\end{widetext}
where $\mathcal{X}_{\rm suff}$ is the state space (without duplicates) associated with the sufficient statistic $\mathcal{X}_{\rm suff} = \{\bm{X}(\bm{x}) | \bm{x} \in \mathcal{X} \}$. Similarly, using Eq.~\eqref{eq:Bayes_2_SM}, we have
\begin{align*}
     \hat{\bm{\gamma}}(\bm{x}) &=  \sum_{y \in \mathcal{Y}} P(y|\bm{x})\bm{\gamma}_{y} = \sum_{y \in \mathcal{Y}} P(y|\bm{x})\bm{\gamma}_{y} P(y|\bm{X}(\bm{x}))\bm{\gamma}_{y}\\
     &= \hat{\bm{\gamma}}\left(\bm{X}(\bm{x})\right) \numberthis,
\end{align*}
and
\begin{align*}\label{eq:SM_der_2}
    \hat{\bm{\gamma}}(\bm{\gamma}) &= \mathbb{E}_{\bm{x} \sim P_{\rm exp}(\bm{x}|\bm{\gamma})}\left[ \hat{\bm{\gamma}}(\bm{x})\right] = \sum_{\bm{x} \in \mathcal{X}} P_{\rm exp}(\bm{x}|\bm{\gamma})\hat{\bm{\gamma}}(\bm{x})\\
    &=  \sum_{\bm{X} \in \mathcal{X}_{\rm suff}} P(\bm{X}|\bm{\gamma}) \hat{\bm{\gamma}}(\bm{X}) = \mathbb{E}_{\bm{X} \sim P(\bm{X}|\bm{\gamma})}\left[  \hat{\bm{\gamma}}(\bm{X})\right]\numberthis.
\end{align*}
The relevant indicators [Eqs. (6), (8), and (10)] of the three phase-classification methods presented in the main text can be straightforwardly computed based on the quantities in Eqs.~\eqref{eq:SM_der_1}, and~\eqref{eq:SM_der_2}, which are expressed solely in terms of the distribution over the sufficient statistic $P(\bm{X}|\bm{\gamma})$ (instead of the full distribution over $\bm{x}$).

In the main text, we considered classical systems with dimensionless Hamiltonians of the form
\begin{equation}\label{eq:exp}
    \mathcal{H}(\bm{x},\bm{\gamma}) = \sum_{i=1}^{d} \gamma_{i} X_{i}(\bm{x}),
\end{equation}
at equilibrium with a large thermal reservoir. In this case, the probability to find the system in state $\bm{x} \in \mathcal{X}$ is given by
\begin{equation}\label{eq:Boltzmann_SM}
   P(\bm{x}| \bm{\gamma}) = e^{-\mathcal{H}(\bm{x},\bm{\gamma})}/Z(\bm{\gamma}).
\end{equation}
Such Boltzmann distributions belong to the exponential family [Eq.~\eqref{eq:SM_exp_fam}] with $h(\bm{x}) = 1$, $A(\bm{\gamma}) = {\rm ln}\,Z(\bm{\gamma})$, and $\eta_{i}(\bm{\gamma}) = \gamma_{i}$. Thus, as a special case of the above, it follows that phase diagrams of such systems can be mapped out given distributions over the sufficient statistic $\{P(\bm{X}|\bm{\gamma})\}_{\bm{\gamma} \in \Gamma}$ (instead of the full distributions over $\bm{x}$).

\section{Data generation}
Here, we provide further details on the data-generation process for the (classical) anisotropic Ising model and the (quantum) cluster-Ising model.
\subsection{Anisotropic Ising model}
Given a set of parameters $\bm{\gamma}=(J_{x}/k_{\rm B}T,J_{y}/k_{\rm B}T)$, we use the Metropolis-Hastings algorithm to sample spin configurations from the corresponding Boltzmann distribution. The lattice is updated by drawing a random spin, which is flipped with probability ${\rm min}(1, e^{-\Delta E/k_{\rm B}T})$, where $\Delta E$ is the energy difference resulting from the considered flip. After a thermalization period of $10^5$ lattice sweeps, we collect $10^{5}$ samples. We treat each quadrant of the phase diagram separately. In each quadrant, we initialize the system in one of the two corresponding ground states and increase $\gamma_{1}$ at constant $\gamma_{2}$. When increasing $\gamma_{1}$, we use the spin configuration from the preceding value as an initial condition for the Markov chain. When increasing $\gamma_{2}$, we reset the system to the corresponding ground state.

\subsection{Cluster Ising model}
Given a set of parameters $\bm{\gamma}=(h_{1}/J,h_{2}/J)$, we compute ground states of the corresponding cluster-Ising Hamiltonian [Eq.~(9) in main text] with $L = 71$ spins using the finite-size density matrix renormalization group (DMRG) with a maximum bond dimension of 150. We utilize the ITensor package~\cite{itensor} in Julia.